\chapter{Introdução}
\label{ch:introducao}

O Capítulo~\ref{ch:introducao} enquadra o problema da priorização
estrutural de requisitos. A Secção~\ref{sec:motivacao} discute a
motivação e os fatores recentes que tornaram esta investigação
oportuna; a Secção~\ref{sec:problema} define formalmente o problema; a
Secção~\ref{sec:objetivos} enuncia os objetivos e as
\emph{research questions} do trabalho; a Secção~\ref{sec:contribuicoes}
lista as contribuições; e a Secção~\ref{sec:estrutura} apresenta a
estrutura do restante documento.

\noindent\textbf{Conceitos-chave:} enquadramento; motivação; problema;
priorização; \emph{research questions}; contribuições.

% ---------------------------------------------------------------------------
\section{Enquadramento e Motivação}
\label{sec:motivacao}

O desenvolvimento de sistemas de software em larga escala enfrenta
desafios crescentes na gestão da complexidade e na organização de
grandes volumes de informação interdependente \parencite{sommerville2007}.
Tradicionalmente, a relevância de elementos dentro de uma rede de
requisitos era determinada por análises textuais, comparações entre
pares ou heurísticas manuais
\parencite{karlsson2007, karlsson1997, berander2005}; a evolução da
\textbf{Teoria dos Grafos} \parencite{newman2010, szwarcfiter2018}
permitiu uma transição para métodos puramente estruturais, capazes de
extrair informação latente diretamente da topologia do sistema.

A motivação principal para este estudo surge da convergência entre
algoritmos de alta performance já em produção na indústria e a
necessidade de rigor analítico na Engenharia de Software. Dois marcos
fundamentais impulsionaram esta investigação:

\begin{itemize}
    \item \textbf{Adopção industrial do SALSA pelo Twitter.}
        Em \emph{Who To Follow}, o sistema de recomendação de
        utilizadores do Twitter, \citet{gupta2013} propõem e descrevem
        uma variante do algoritmo \textbf{SALSA}~\parencite{lempel2001}
        para sugerir conexões entre utilizadores. Esta adopção em
        produção, sobre um grafo de larga escala, demonstra a
        viabilidade de usar passeios aleatórios em grafos para
        determinar prestígio e relevância de forma escalável.

    \item \textbf{Aplicação do PageRank à priorização de requisitos.}
        A descoberta de pesquisas que aplicam o
        \textbf{PageRank}~\parencite{page1999} à priorização de
        requisitos de software
        \parencite{silva2018, jin2009, li2009} sustenta a hipótese de
        que a estrutura de dependências de um sistema contém
        informação latente suficiente para ordenar requisitos de forma
        mais consistente do que técnicas predominantemente subjetivas
        \parencite{firesmith2004, achimugu2014}.
\end{itemize}

Dessa forma, o trabalho enquadra-se na busca por modelos matemáticos
que minimizem o viés humano em decisões críticas de engenharia,
utilizando propriedades topológicas para extrair inteligência de grafos
de software.

% ---------------------------------------------------------------------------
\section{Problema}
\label{sec:problema}

O \textbf{ranqueamento de requisitos} --- termo correntemente
designado na literatura por \textbf{priorização de requisitos}
\parencite{karlsson2002, firesmith2004, silva2018, bukhsh2020} ---
consiste em produzir uma ordem total (ou parcial) sobre o conjunto de
requisitos $R = \{r_1, r_2, \ldots, r_n\}$ de um sistema, a partir de
critérios que reflitam a importância relativa de cada requisito para a
entrega de valor, a estabilidade arquitetural ou o risco do projeto.
Formalmente, dado o conjunto $R$ e uma relação de dependência
$D \subseteq R \times R$, em que $(r_i, r_j) \in D$ se $r_i$ depende
de $r_j$, a tarefa de priorização procura uma função
$\pi: R \rightarrow \mathbb{R}$ que atribua a cada requisito um
\emph{score} compatível com a sua centralidade na rede $(R, D)$.

\begin{notebox}[Nota ontológica]
    No presente trabalho, ``requisito'' é uma \textbf{agregação de
    responsabilidades implementadas no código-fonte}, identificada
    via um mapeamento explícito \code{func\_to\_req} (Secção
    \ref{ch:metodologia}). Não corresponde, portanto, a um requisito
    funcional clássico no sentido de \citet{sommerville2007}, mas
    serve como \emph{proxy} de \emph{domain feature} para efeitos de
    ranqueamento estrutural. Esta opção é discutida em detalhe na
    Secção~\ref{sec:posicionamento} (posicionamento face à literatura)
    e nas ameaças à validade (Secção~\ref{sec:ameacas}).
\end{notebox}

As técnicas tradicionais de priorização --- Comparação em Pares
\parencite{karlsson2007}, Votação Acumulativa
\parencite{leffingwell2003, ahl2005}, AHP (\emph{Analytic Hierarchy
Process}) \parencite{saaty1980, karlsson1997} e PERT/CPM
\parencite{kerzner2009} --- sofrem de problemas intrínsecos de
subjetividade e esforço \parencite{firesmith2004}. A inexperiência dos
\emph{stakeholders}, a divergência na interpretação das escalas de
prioridade e o foco excessivo num único ponto de vista resultam em
prioridades inconsistentes e em ``posições inválidas''
\parencite{silva2018}. Em particular, a Votação Acumulativa, por
depender excessivamente da decisão direta do \emph{stakeholder}, não
resolve a duplicidade envolvida em requisitos mutuamente dependentes
\parencite{silva2018}.

Este trabalho justifica-se, portanto, ao investigar e comparar a
eficácia do \textbf{PageRank}, \textbf{HITS} e \textbf{SALSA} na
priorização estrutural de requisitos, sobre grafos extraídos
automaticamente do código-fonte, oferecendo um modelo analítico que
minimize o esforço dos \emph{stakeholders} e que utilize integralmente
a rastreabilidade de requisitos disponível
\parencite{gotel1994, rasiman2022}.

% ---------------------------------------------------------------------------
\section{Objetivos e \emph{Research Questions}}
\label{sec:objetivos}

O objetivo desta dissertação é \textbf{investigar e comparar a
eficácia} dos algoritmos PageRank, HITS e SALSA na priorização
estrutural de requisitos de software, fornecendo um modelo analítico
que utilize a rastreabilidade entre requisitos para garantir
consistência em sistemas de larga escala. Em concreto, pretende-se:
(i) \textbf{formalizar} a aplicação de processos estocásticos e de
abordagens de Teoria dos Grafos à modelação de interdependências entre
requisitos e à determinação de fluxos de influência;
(ii) \textbf{desenvolver} um protótipo (ReqGraph) que extraia, de
forma totalmente automatizada e reprodutível, o grafo de requisitos a
partir do código-fonte de um projeto Python; (iii) \textbf{avaliar
empiricamente}, por meio de experimentação sobre projetos reais de
código aberto, o comportamento comparativo dos três algoritmos,
identificando convergências e divergências de ranking e discutindo-as
à luz da arquitetura conhecida dos sistemas analisados.

Estes objetivos materializam-se nas seguintes \textit{research
questions} (RQs), formuladas de forma verificável e respondidas no
Capítulo~\ref{ch:resultados}:

\begin{description}
    \item[RQ1.] Em que medida os rankings produzidos por PageRank,
        HITS e SALSA \emph{concordam} entre si quando aplicados a
        grafos de requisitos extraídos automaticamente por análise
        estática (AST)? A concordância é medida pelo coeficiente de
        Spearman ($\rho$) e Kendall ($\tau_b$) entre cada par de
        algoritmos.
    \item[RQ2.] Como varia essa concordância em função das
        propriedades estruturais do grafo (densidade, presença de
        ciclos, simetria de graus)?
    \item[RQ3.] Em que medida os \emph{top-$k$ requisitos} apontados
        pelos algoritmos correspondem a elementos conhecidos como
        arquiteturalmente centrais nos projetos analisados (\ie\
        qual a coerência arquitetural dos rankings produzidos)?
\end{description}

A formulação destas RQs permite distinguir o que esta investigação
\emph{demonstra} (acordo/desacordo entre algoritmos sobre os mesmos
dados) do que \emph{não} demonstra (validade preditiva face a juízos
humanos --- ver Secção~\ref{sec:ameacas} e
Secção~\ref{sec:trabalho-futuro}).

% ---------------------------------------------------------------------------
\section{Contribuições}
\label{sec:contribuicoes}

As principais contribuições desta dissertação são:

\begin{enumerate}
    \item \textbf{Pipeline reprodutível para extração de grafos de
        requisitos.} Implementação do protótipo \emph{ReqGraph}, que
        automatiza a transformação \emph{Código-fonte
        \(\rightarrow\) AST \(\rightarrow\) Call Graph
        \(\rightarrow\) Grafo de Requisitos} sobre projetos Python,
        recorrendo exclusivamente a análise estática e a artefactos
        canónicos da comunidade científica (NetworkX, Graphviz).
    \item \textbf{Implementação aberta do SALSA.} Implementação manual
        e documentada do algoritmo SALSA \parencite{lempel2001}, uma
        vez que o NetworkX não a disponibiliza nativamente, incluindo
        o tratamento explícito de vértices absorventes por
        redistribuição uniforme.
    \item \textbf{Comparação empírica quantitativa entre PageRank,
        HITS e SALSA} sobre grafos reais de requisitos derivados de
        três projetos \emph{open source} representativos de regimes
        crescentes de acoplamento arquitetural (CPython \emph{stdlib},
        Flask e \emph{scikit-learn}). A comparação utiliza métricas
        formais de correlação de ranks (Spearman $\rho$, Kendall
        $\tau_b$) entre cada par de algoritmos.
    \item \textbf{Mecanismo de redução do custo de mapeamento.}
        Formalização de um \emph{prompt} estruturado
        (\code{llm\_prompt\_mapping.md}) que permite gerar o
        dicionário \code{func\_to\_req} via modelos de linguagem de
        grande escala, mitigando a principal fricção operacional da
        abordagem.
\end{enumerate}

O código-fonte completo, os mapeamentos \code{func\_to\_req}
utilizados e os artefactos gerados nas experiências encontram-se
publicamente disponíveis em
\url{https://github.com/tales96323/Tales_Tese_ulusofona}.

% ---------------------------------------------------------------------------
\section{Estrutura do Documento}
\label{sec:estrutura}

O restante documento está organizado da seguinte forma. O
\textbf{Capítulo~\ref{ch:fundamentos}} apresenta os conceitos teóricos
relevantes --- Teoria dos Grafos, algoritmos de ranqueamento e
fundamentos da Engenharia de Requisitos --- e contextualiza o
problema, posicionando esta dissertação em relação ao trabalho de
\citet{silva2018}, do qual constitui uma evolução metodológica. O
\textbf{Capítulo~\ref{ch:metodologia}} descreve a metodologia
adotada, a arquitetura do protótipo ReqGraph, o protocolo experimental
e o protocolo de mapeamento \code{func\_to\_req}. O
\textbf{Capítulo~\ref{ch:resultados}} apresenta e discute os
resultados obtidos nos três casos de estudo, recorrendo a métricas
formais de correlação de ranks e a uma análise dedicada de
\emph{ameaças à validade} segundo o referencial de
\citet{wohlin2012}. O \textbf{Capítulo~\ref{ch:conclusao}} conclui o
trabalho, sintetiza as contribuições e identifica linhas de trabalho
futuro.
